% TODO MATEJ: remove versions
\subsection{Background}

% What is the background of the company for which you carried out this project? Not applicable if the project was not conducted for a company. \label{q} \newline
The project was carried out in company Result d.o.o. They are a software development company, but also focus on data, artificial intelligence, and business solutions. They work on many projects, some of which are european. One such is eEarly --- an open platform for analysation and gathering of health data\cite{EEarly}. It was also the umbrella for this bachelor work.

% Why would the company want this to be researched or developed? If not for a company: what is the motivation behind the bachelor project? \label{q} \newline
European Union has many regulations and laws when it comes to data, and there are even more for each individual country, organisation,\ldots This introduces challenges for projects like eEarly that utilize artificial intelligence and deal with sensitive medical data. From here we can easily see the motivation for my bachelor project: To find the bridge between sensitive data, artificial intelligence, and a lake of regulations.


\subsection{Problem statement}

% What is the problem statement? \label{q} \newline
% TODO Q&A MANJA: "each hospital keeps its data isolated" source
In an ideal situation, an artificial intelligence model could train on all data without restrictions. In reality, besides strict European Union laws, each hospital keeps its data isolated. This reveals the main issue: An artificial intelligence model can only train on one hospital. Consequences of this are bad model performance, predictions, and ultimately uselessness of the new technology. From this follows the proposal, which is not to loosen the law, but to explore techniques like federated learning\cite{FederatedLearning2026} to satisfy both the regulations and the artificial intelligence's need for data.


\subsection{Intention}

% What is the company’s intention? If not for a company: how will the product be used in the long term (independently of a company)? \label{q} \newline
The main intention of the project is to develop a federated learning proof of concept for the company, which can be later on used as a starting point for other projects and might be fully implemented, possibly in a healthcare environment, or be a new privacy-oriented artificial intelligence solution that the company could sell.


\subsection{Scope within the bachelor project (MVP – need to have)}

% What is the scope of the project within this bachelor project? \label{q} \newline
The scope includes terminology described in section~\ref{sec:preliminary-research}. ``Must have'' features describe the minimal viable product (MVP); without these, the project is not deemed finished, unless blockers are clearly stated.

\begin{table}[htb]\centering
\caption{Features with their importance. At the top are the most important ones.}\label{tab:scope}
 \csvreader[
  tabular = | p{0.2\textwidth} | p{0.75\textwidth} |,
  table head = \hline \textbf{Importance} & \textbf{Feature}\\\hline,
  late after line = \\\hline,
  separator = semicolon,
  caption = {My table}
 ]{media/scope.csv}%
 {Importance=\importance, Feature=\feature}{%
  \importance & \feature
 }
\end{table}


\subsection{Ultimate scope (want to have)}

%What is the ultimate scope? \label{q} \newline

Ultimate scope of the project is defined with features of ``Should have'' and ``Could have'' importance in the table~\ref{tab:scope}. As you can see, as importance lowers, features become closer to those of realistic setting in which federated learning would have been deployed.